\counterwithin{figure}{section}
\counterwithin{table}{section}
\counterwithin{equation}{section}
\chapter{Introducción}

Las correas transportadoras constituyen un componente crítico en la minería y en diversas industrias intensivas en el transporte de materiales a granel, al permitir el movimiento continuo y eficiente de grandes volúmenes de carga sobre largas distancias. Su operación sostenida se encuentra estrechamente vinculada a la productividad global de la planta, de modo que fallas inesperadas en la correa o en sus componentes asociados se traducen directamente en detenciones no programadas, sobrecostos de mantenimiento y potenciales riesgos para la seguridad de las personas y equipos. En este contexto, asegurar la disponibilidad y confiabilidad de las correas transportadoras es un requisito fundamental para la operación industrial moderna. \\

Durante su vida útil, las correas están expuestas a condiciones de servicio severas que favorecen la aparición de diversos modos de daño, tales como desgaste abrasivo, grietas o rasgaduras, roturas de empalme y fallas asociadas a desalineamientos prolongados. Estos defectos pueden evolucionar desde deterioros superficiales hasta fallas catastróficas, comprometiendo la integridad estructural de la correa y de la línea de transporte \cite{khanna2022types}. La gran longitud de los sistemas transportadores y la variabilidad de las condiciones ambientales dificultan su inspección exhaustiva, incrementando la probabilidad de que daños incipientes no sean detectados oportunamente. \\

Tradicionalmente, el monitoreo del estado de las correas transportadoras se ha basado en rondas de inspección visual realizadas por personal especializado, apoyadas ocasionalmente por dispositivos de medición puntual. Si bien este enfoque ha sido la práctica estándar en la industria, presenta limitaciones importantes: la cobertura espacial y temporal de las inspecciones es restringida, la detección de fallas depende fuertemente de la experiencia del inspector y, en muchos casos, es necesario detener o intervenir la correa para realizar revisiones detalladas, afectando la continuidad operacional. Estas restricciones han motivado el desarrollo de estrategias de monitoreo automatizado que reduzcan la carga operativa y mejoren la detección temprana de daños \cite{angulo2024deteccion}. \\

En paralelo, los avances en visión computacional y aprendizaje profundo (\textit{Deep Learning}) han permitido desarrollar modelos capaces de identificar y localizar objetos o patrones específicos directamente a partir de imágenes, sin requerir una etapa manual de ingeniería de características. En particular, los detectores de una sola etapa (\textit{one-stage detectors}), como las familias YOLO y SSD, han mostrado un compromiso favorable entre precisión y velocidad de inferencia en escenarios de detección en tiempo casi real, lo que los convierte en candidatos naturales para su aplicación en el monitoreo continuo de correas transportadoras \cite{zhang2021deep,zhang2021one_stage}. Adicionalmente, trabajos recientes han demostrado la viabilidad de utilizar arquitecturas basadas en YOLO para la detección de fallas o condiciones anómalas en sistemas transportadores, reforzando el potencial de estas técnicas en entornos industriales \cite{angulo2024deteccion}. \\

Considerando este contexto, la presente memoria aborda el problema de la identificación automática de fallas en correas transportadoras mediante modelos de detección de objetos basados en aprendizaje profundo, abarcando tanto detectores de una sola etapa como arquitecturas recientes basadas en \textit{transformers}. El trabajo se centra en la implementación, entrenamiento y evaluación comparativa de estos enfoques sobre un conjunto de datos específico de defectos en correas, con el objetivo de analizar su capacidad para detectar distintos tipos de daño y su adecuación para ser integrados en sistemas de monitoreo y mantenimiento predictivo en un entorno cercano al operativo, según la eficiencia y exigencia de cada modelo. \\





%%%%%%%%%%%%%%%%%%%%%%%%%%%%%%%%%%%%%%%%%%%%
\section{Objetivo General}

Implementar y evaluar algoritmos
de detección de objetos, basados en
inteligencia artificial, para la identificación automática de
fallas en correas transportadoras.



\section{Objetivos Específicos}


Para lograr el objetivo general en este futuro trabajo de memoria de título, se propone cumplir con los siguientes
objetivos específicos:

\begin{enumerate}
    \item Implementar modelos de detección de objetos basados en redes convolucionales, tales como YOLO (You Only Look Once) y sus variantes, y SSD (Single Shot Multibox Detector).

    \item Implementar modelos de detección de objetos basados en redes neuronales con arquitectura tipo transformer, tales como DETR (DEtection TRansformer) y DINO.

    \item Entrenar algoritmos de detección de objetos utilizando bases de datos con registros de fallas en correas transportadoras.

    \item Comparar el desempeño de los modelos en la detección de distintos tipos de fallas mediante el uso de métricas de evaluación cuantitativas.

    
\end{enumerate}


\section{Alcances}

\begin{itemize}
    \item El entrenamiento de los algoritmos de detección se realizará exclusivamente con bases de datos públicas disponibles en la web que contengan registros de fallas en correas transportadoras.

    \item Considerar el uso de datos sin etiquetar, los cuales serán procesados mediante estrategias de etiquetado manual o mediante técnicas no supervisadas y/o auto-supervisadas, según su aplicabilidad.

    \item Evaluar el rendimiento de los modelos seleccionados (YOLO, SSD, DETR y DINO) en la detección y clasificación de fallas, utilizando una misma base de datos con partición en conjuntos de entrenamiento, validación y prueba. Posteriormente se evaluarán utilizando métricas estándar de evaluación.
    
\end{itemize}